\documentclass[12pt,a4paper,oneside]{report}
\usepackage[utf8]{inputenc}
\usepackage[italian]{babel}
\usepackage{graphicx}
\usepackage{hyperref}
\usepackage{geometry}
\usepackage{listings}
\usepackage{xcolor}
\usepackage{titlesec}
\usepackage{float}
\usepackage{fancyhdr}
\usepackage{microtype}
\usepackage{tikz}
\usepackage{booktabs}
\usepackage{longtable}
\usepackage{array}
\usepackage{tabularx}
\usepackage{multirow}
\usepackage{enumitem}
\usepackage{amsmath}
\usepackage{amssymb}
\usepackage{tcolorbox}
\tcbuselibrary{skins,breakable}
\usetikzlibrary{positioning, shapes.geometric, arrows.meta, fit, backgrounds, calc}

\geometry{a4paper, top=2.8cm, bottom=2.8cm, left=3cm, right=3cm}
\setlength{\headheight}{13.6pt}
\addtolength{\topmargin}{-1.6pt}
\hypersetup{
    colorlinks=true,
    linkcolor=blue!70!black,
    filecolor=magenta,
    urlcolor=cyan!60!black,
    pdftitle={Relazione Progetto Cloud - DocuMatch},
    pdfauthor={Manuel Amodio},
    pdfsubject={Sistemi Distribuiti e Cloud Computing}
}

\definecolor{codebg}{RGB}{245,247,250}
\definecolor{codeframe}{RGB}{200,210,220}
\definecolor{commentgreen}{RGB}{34,139,34}
\definecolor{keywordblue}{RGB}{0,0,180}
\definecolor{stringred}{RGB}{180,0,0}
\definecolor{promptbg}{RGB}{240,248,255}
\definecolor{promptframe}{RGB}{100,149,237}
\definecolor{warningbg}{RGB}{255,248,230}
\definecolor{warningframe}{RGB}{255,165,0}

\lstdefinestyle{mypython}{
    backgroundcolor=\color{codebg},
    commentstyle=\color{commentgreen}\itshape,
    keywordstyle=\color{keywordblue}\bfseries,
    stringstyle=\color{stringred},
    basicstyle=\ttfamily\footnotesize,
    breakatwhitespace=false,
    breaklines=true,
    captionpos=b,
    keepspaces=true,
    numbers=left,
    numbersep=8pt,
    numberstyle=\tiny\color{gray},
    showspaces=false,
    showstringspaces=false,
    showtabs=false,
    tabsize=2,
    frame=single,
    rulecolor=\color{codeframe},
    language=Python
}
\lstdefinestyle{mybash}{
    backgroundcolor=\color{codebg},
    basicstyle=\ttfamily\footnotesize,
    breaklines=true,
    frame=single,
    rulecolor=\color{codeframe},
    numbers=none
}
\lstset{style=mypython}

\newtcolorbox{promptbox}[1][]{
    colback=promptbg,
    colframe=promptframe,
    fonttitle=\bfseries\sffamily,
    title=Prompt AI Generativa,
    breakable,
    #1
}
\newtcolorbox{notabox}[1][]{
    colback=warningbg,
    colframe=warningframe,
    fonttitle=\bfseries\sffamily,
    title=Nota Tecnica,
    breakable,
    #1
}
\pagestyle{plain}
\begin{document}

\begin{titlepage}
    \centering
    \vspace*{0.8cm}
    {\scshape\LARGE Universit\`{a} della Calabria \par}
    \vspace{0.4cm}
    {\scshape\Large Dipartimento di Ingegneria Informatica, Modellistica,\\
    Elettronica e Sistemistica (DIMES)\par}
    \vspace{0.3cm}
    {\large Corso di Laurea Magistrale in Ingegneria Informatica(Cybersecurity)\par}
    \vspace{0.2cm}
    {\large Corso di Sistemi Distribuiti e Cloud Computing --- A.A.\ 2025/2026\par}
    \vspace{1.2cm}
    \rule{\linewidth}{1pt}
    \vspace{0.4cm}
    {\huge\bfseries DocuMatch\par}
    \vspace{0.3cm}
    {\LARGE\bfseries Relazione di Progetto\par}
    \vspace{0.2cm}
    {\Large\textit{Sistema Cloud-Native ed Intelligente per la Gestione,\\
    Ricerca Semantica Vettoriale ed Estrazione OCR\\
    di schede Documentali Comunali tramite AI Generativa}\par}
    \vspace{0.4cm}
    \rule{\linewidth}{1pt}
    \vspace{1.2cm}
    \begin{minipage}[t]{0.48\textwidth}
        \begin{flushleft}\large
            \emph{Candidato:}\\[4pt]
            \textbf{Manuel \textsc{Amodio}}\\
            Matricola: \textbf{276571}\\
        \end{flushleft}
    \end{minipage}
    \hfill
    \begin{minipage}[t]{0.48\textwidth}
        \begin{flushright}\large
            \emph{Docenti:}\\[4pt]
            \textbf{Prof. Domenico \textsc{Talia}}\\
            \textbf{Prof. Loris \textsc{Belcastro}}\\
        \end{flushright}
    \end{minipage}
    \vfill
    {\large Piattaforma Cloud: \textbf{Microsoft Azure}\par}
    \vspace{0.2cm}
    {\large Anno Accademico 2025/2026\par}
\end{titlepage}

\tableofcontents
\newpage

\chapter{Introduzione e Requisiti del Progetto}

\section{Contesto e Motivazione}

La digitalizzazione dei documenti amministrativi comunali \`{e} una delle sfide pi\`{u}
concrete e diffuse nella Pubblica Amministrazione italiana. I documenti amministrativi, restano spesso disponibili soltanto come
scansioni non ricercabili o come fascicoli cartacei. Questo rende estremamente onerosa
la consultazione per i funzionari e quasi impossibile per i comuni cittadini.

Il progetto \textbf{DocuMatch} nasce per rispondere a questa esigenza reale, sfruttando
le pi\`{u} moderne tecnologie di \textit{Cloud Computing}, \textit{Intelligenza Artificiale}
e \textit{Ricerca Vettoriale Semantica} per costruire un sistema documentale istituzionale
completo, scalabile e accessibile.

L'applicativo supera i limiti della classica ricerca testuale (keyword-based) integrando
strumenti di AI che comprendono il significato della query dell'utente indipendentemente
dalle parole esatte utilizzate.

\section{Requisiti Funzionali}

Tutti gli obiettivi definiti dalla traccia assegnata sono stati pienamente raggiunti:

\begin{enumerate}[leftmargin=2cm]
    \item \textbf{Architettura Cloud-Native su Microsoft Azure:} esecuzione su servizi
          gestiti Azure (Container Apps, Blob Storage, Document Intelligence, OpenAI).
    \item \textbf{OCR Automatico:} estrazione del testo dalle copertine tramite Azure AI
          Document Intelligence.
    \item \textbf{Ricerca Semantica Vettoriale:} indicizzazione tramite embedding
          \texttt{paraphrase-multilingual-MiniLM-L12-v2} e cosine similarity.
    \item \textbf{Caricamento Massivo:} supporto CSV/JSON con calcolo
          automatico embedding per ogni record.
    \item \textbf{AI Generativa:} Azure OpenAI per spiegazioni e chatbot RAG.
    \item \textbf{Notifiche Schedulate:} APScheduler invia email agli operatori per
          atti in scadenza a 15, 7 e 1 giorno.
    \item \textbf{Autenticazione e Sicurezza:} token HMAC-SHA256, bcrypt, reset password.
    \item \textbf{Containerizzazione:} Docker + Docker Compose + Azure Container Apps.
    \item \textbf{Interfaccia Istituzionale:} SPA React con UI custom esportata tramite FigmaMake.
\end{enumerate}

\section{Requisiti Non Funzionali}

\begin{itemize}
    \item \textbf{Portabilit\`{a}:} pattern \textit{Dual-Mode} --- sviluppo offline senza utilizzo di servizi cloud,
          deploy cloud senza modifiche al codice.
    \item \textbf{Scalabilit\`{a}:} Azure Container Apps con auto-scaling su Kubernetes.
    \item \textbf{Manutenibilit\`{a}:} codice SOLID (DIP), struttura a 3 strati.
    \item \textbf{Sicurezza:} CORS per dominio, bcrypt, token con scadenza, anti-timing.
    \item \textbf{Osservabilit\`{a}:} healthcheck Docker, endpoint root JSON per load balancer.
\end{itemize}

\chapter{Architettura del Sistema}

\section{Panoramica Generale}

DocuMatch adotta un'architettura \textbf{Client-Server a tre strati} con Separation of Concerns:

\begin{enumerate}
    \item \textbf{Presentation Layer (Frontend):} SPA React 19 + TypeScript + Vite,
          servita da container Nginx con reverse proxy.
    \item \textbf{Business Logic Layer (Backend):} API REST FastAPI (Python 3.9),
          orchestratore di tutti i servizi (NLP, OCR, storage, email, AI).
    \item \textbf{Data \& Storage Layer:} SQLite (sviluppo) / PostgreSQL (produzione)
          tramite SQLAlchemy ORM; Azure Blob Storage per file binari.
\end{enumerate}

\section{Schema Architetturale}

\begin{figure}[H]
    \centering
    \begin{tikzpicture}[
        node distance=1.8cm and 1.5cm,
        box/.style={rectangle, draw=black!70, thick, fill=white,
                    text width=2.5cm, align=center, rounded corners=4pt,
                    minimum height=1.1cm, font=\small},
        cloudbox/.style={ellipse, draw=blue!60, thick, fill=blue!5,
                         text width=2.2cm, align=center, font=\small},
        orangebox/.style={rectangle, draw=orange!70, thick, fill=orange!8,
                          text width=2.5cm, align=center, rounded corners=4pt,
                          minimum height=1.1cm, font=\small},
        db/.style={cylinder, draw=gray!70, thick, aspect=0.25,
                   fill=gray!10, text width=1.9cm, align=center,
                   minimum height=1.3cm, shape border rotate=90, font=\small},
        arrow/.style={-{Stealth[length=3mm]}, thick, draw=gray!70}
    ]
    \node[box, fill=blue!10] (browser) {Browser\\{\scriptsize React 19 / TS}};
    \node[box, fill=green!10, below=of browser] (nginx) {Nginx\\{\scriptsize Rev. Proxy}};
    \node[box, fill=green!20, below=of nginx] (fastapi) {FastAPI\\{\scriptsize Python 3.9}};
    
    \node[db, left=of fastapi] (db) {SQLite/\\PostgreSQL};
    \node[orangebox, right=of fastapi] (openai) {Azure OpenAI\\{\scriptsize GPT 5.4 Mini}};
    
    \node[cloudbox, below left=of fastapi] (blob) {Azure Blob\\{\scriptsize Storage}};
    \node[cloudbox, fill=purple!5, below right=of fastapi] (ocr) {Azure Doc.\\{\scriptsize Intelligence}};
    
    \node[box, fill=yellow!10, left=of nginx] (sched) {APScheduler\\{\scriptsize Cron 08:00}};
    
    \draw[arrow, <->] (browser) -- node[right,font=\scriptsize]{HTTP/REST} (nginx);
    \draw[arrow, <->] (nginx) -- node[right,font=\scriptsize]{Proxy} (fastapi);
    \draw[arrow, <->] (fastapi) -- node[above,font=\scriptsize]{SQLAlchemy} (db);
    \draw[arrow, <->] (fastapi) -- node[above,font=\scriptsize]{REST} (openai);
    \draw[arrow, <->] (fastapi) -- node[above left,font=\scriptsize, sloped]{SDK} (blob);
    \draw[arrow, <->] (fastapi) -- node[above right,font=\scriptsize, sloped]{SDK} (ocr);
    \draw[arrow, ->] (sched) -- node[above left,font=\scriptsize, sloped]{trigger} (fastapi);
    \coordinate (box_bottom) at ([yshift=-3cm]fastapi.south);
    \coordinate (box_top) at ([yshift=2.5cm]nginx.north);
    \begin{scope}[on background layer]
        \node[draw=gray!50, dashed, thick, fill=gray!3,
              fit=(nginx)(fastapi)(db)(openai)(sched)(box_bottom)(box_top),
              inner sep=0.6cm,
              label={[anchor=south,yshift=-0.1cm]south:
                \footnotesize\textbf{Azure Container Apps (Backend)}}]
              (backendenv) {};
    \end{scope}
    \end{tikzpicture}
    \caption{Schema architetturale di DocuMatch con tutti i servizi Azure}
    \label{fig:arch}
\end{figure}

Il flusso di una ricerca semantica:
\begin{enumerate}
    \item Il browser invia \texttt{GET /api/v1/documents/search?query\_testo=lavori+stradali}.
    \item Nginx la inoltrata al container backend FastAPI.
    \item FastAPI genera l'embedding della query con SentenceTransformers.
    \item Calcola la cosine similarity rispetto a tutti i documenti nel DB.
    \item La risposta JSON torna al frontend con card e punteggi di rilevanza.
\end{enumerate}
\vspace{4cm}
\section{Stack Tecnologico}

\begin{longtable}{>{\bfseries}p{3.2cm} p{4.2cm} p{5cm}}
    \toprule
    \textbf{Componente} & \textbf{Libreria/Versione} & \textbf{Motivazione} \\
    \midrule
    \endhead
    Framework API & FastAPI 0.110.0 & ASGI, Swagger auto-generato, async nativo \\
    Server ASGI & Uvicorn 0.28.0 & Ad alte prestazioni per FastAPI \\
    ORM Database & SQLAlchemy 2.0.28 & Astrazione SQLite/PostgreSQL trasparente \\
    Validazione & Pydantic 2.6.4 & Schema dati con validazione automatica \\
    NLP / NER & spaCy 3.7.4 & Riconoscimento entit\`{a} in italiano \\
    Embeddings & SentenceTransformers 2.5.1 & Vettori semantici multilingua a 384 dim \\
    Calcolo vettoriale & NumPy 1.26.4 & Cosine similarity e operazioni vettoriali \\
    Importazione dati & Pandas 2.2.1 & Parsing CSV con separatori ambigui \\
    Hashing password & Passlib[bcrypt] 1.7.4 & Key stretching resistente a GPU \\
    Token auth & HMAC-SHA256 (stdlib) & Stateless con firma e scadenza \\
    Scheduler & APScheduler 3.10.4 & Cron job in background \\
    Azure OCR & azure-ai-formrecognizer 3.3.2 & Document Intelligence SDK \\
    Azure Storage & azure-storage-blob 12.19.1 & Blob Storage SDK \\
    Azure AI & openai 1.14.1 & Client Azure OpenAI e OpenAI API \\
    DB Produzione & psycopg2-binary 2.9.9 & Driver PostgreSQL \\
    Build tool frontend & Vite & HMR istantaneo, ESBuild \\
    UI framework & React 19 + TypeScript & Componenti dichiarativi, design to code \\
    \bottomrule
    \caption{Stack tecnologico completo di DocuMatch}
\end{longtable}

\section{Alternative Valutate e Scartate}

\begin{description}
    \item[VM (IaaS) vs Container Apps (PaaS):] La VM richiede gestione del Guest OS.
        Container Apps (su Kubernetes gestito) astrae l'infrastruttura e offre
        auto-scaling, networking e TLS automatici.
    \item[SQL LIKE vs Ricerca Semantica:] \texttt{LIKE '\%lavori stradali\%'} non trova
        ``rifacimento manto stradale''. La cosine similarity su vettori semantici
        cattura la similariit\`{a} concettuale.
    \item[Tesseract vs Azure AI Document Intelligence:] Tesseract ha performance
        inferiori su scansioni reali di bassa qualit\`{a}. Il servizio cloud \`{e}
        ottimizzato per documenti burocratici e atti ufficiali.
    \item[GPT 5.6 Sol vs GPT 5.4 Mini:] Inizialmente si era valutato l'utilizzo del modello \textbf{GPT 5.6 Sol} (estremamente preciso), ma i suoi tempi di risposta prolungati sono risultati incompatibili con la necessit\`{a} di interazione rapida e fluida richiesta dal cliente, che non vuole attendere. Si \`{e} scelto quindi \textbf{GPT 5.4 Mini}: pur essendo intrinsecamente meno preciso, l'adozione di un \textit{prompt engineering} aggiuntivo ha permesso di colmare il divario qualitativo, raggiungendo gli stessi risultati in maniera decisamente pi\`{u} veloce.
    \item[Azure Communication Services vs Gmail SMTP:] A causa di rigide limitazioni imposte da Azure per l'abilitazione del servizio di posta elettronica sui nuovi account, si \`{e} optato per l'affidabile e immediato server SMTP di Gmail per la gestione delle email transazionali.
    \item[SQLite in produzione vs PostgreSQL:] SQLite non \`{e} adatto al deployment
        multi-processo. L'astrazione SQLAlchemy rende il cambio trasparente.
\end{description}

\chapter{Layer di Persistenza e Modello dei Dati}

\section{Architettura Dual-Mode del Database}

Il backend non dipende da un database specifico ma dall'astrazione\textbf{ SQLAlchemy}.
La connection string viene letta a runtime da \texttt{DATABASE\_URL}:
In sviluppo viene usato SQLite (file locale), in produzione PostgreSQL su Azure.
Il cambio \`{e} completamente trasparente per il codice applicativo.

\section{Modelli SQLAlchemy}

\subsection{Modello \texttt{Document}}

Entit\`{a} centrale. Le liste (uffici, firmatari, frazioni) sono serializzate come JSON string
poich\'{e} SQLite non supporta array nativi:

\subsection{Modello \texttt{User}}

Gestisce l'autenticazione degli operatori e l'eventuale recupero password. Le caratteristiche principali includono:
\begin{itemize}
    \item \textbf{ID UUID (String 36)}: Utilizzato al posto di ID sequenziali per prevenire attacchi di tipo Insecure Direct Object Reference (IDOR).
    \item \textbf{Hashing bcrypt}: La password non viene mai salvata in chiaro ma protetta da algoritmi crittografici \textit{salt-based}.
    \item \textbf{Flow di Reset Password}: Comprende i campi \texttt{reset\_token} (monouso) e \texttt{reset\_token\_expiry} (validit\`{a} 15 minuti) per consentire un ripristino sicuro delle credenziali in caso di smarrimento.
\end{itemize}

\subsection{Pannello di Amministrazione e Configurazione}

Il pannello admin permette la gestione centralizzata delle seguenti entit\`{a} e sezioni:
\begin{itemize}
    \item \textbf{Operatori}: gestione degli utenti del sistema (creazione account, reset password, permessi).
    \item \textbf{Registro Attivit\`{a}}: monitoraggio delle operazioni (Audit Log) per garantire la tracciabilit\`{a} e la sicurezza.
    \item \textbf{Uffici}: uffici comunali (Ufficio Tecnico, Servizi Sociali, \ldots)
    \item \textbf{Tipologie}: categorie documentali (Delibera di Giunta, Ordinanza, \ldots)
    \item \textbf{Firmatari}: funzionari abilitati alla firma degli atti.
    \item \textbf{Frazioni}: zone geografiche del territorio comunale.
\end{itemize}

\chapter{Servizi Cloud Microsoft Azure}

\section{Pattern Dual-Mode e Dependency Inversion}

Ogni servizio cloud adotta il \textbf{design pattern Strategy} con DIP (SOLID):
\begin{enumerate}
    \item Classe astratta (contratto).
    \item Implementazione cloud con API Azure reali.
    \item Implementazione mock funzionale per sviluppo offline.
\end{enumerate}
La selezione avviene a runtime: se le variabili d'ambiente Azure sono impostate si usa
il servizio cloud; altrimenti si attiva automaticamente il mock locale.

\section{Azure AI Document Intelligence (OCR)}

\textbf{Motivazione:} garantisce alta accuratezza su scansioni di qualit\`{a} variabile,
gestisce layout complessi e tabelle degli atti ufficiali. A differenza di Tesseract, il
servizio cloud \`{e} ottimizzato per documenti burocratici reali. Si usa il modello
pre-addestrato \texttt{prebuilt-document}.

\section{Azure Blob Storage}

\textbf{Motivazione:} Object Storage per immagini delle copertine. Vantaggi rispetto al disco locale:
\begin{itemize}
    \item \textbf{Persistenza stateless:} il container pu\`{o} essere riavviato/replicato senza perdere file.
    \item \textbf{URL pubblici permanenti:} ogni blob ha un URL diretto senza endpoint di download.
    \item \textbf{Ridondanza geografica (LRS/GRS):} replica automatica per alta disponibilit\`{a}.
\end{itemize}

\section{Azure OpenAI Service (Chatbot RAG ed Estrazione Metadati)}

Integra il modello \texttt{GPT 5.4 Mini} per:
\begin{enumerate}
    \item \textbf{Estrazione Metadati OCR:} trasformazione del testo OCR destrutturato in JSON strutturato.
    \item \textbf{Spiegazioni semantiche:} linguaggio naturale per i risultati dei documenti correlati.
    \item \textbf{Chatbot RAG:} interagisce con l'archivio testuale limitando le allucinazioni.
\end{enumerate}

\subsection{Architettura del Chatbot RAG}

Il chatbot segue il pattern \textbf{RAG (Retrieval-Augmented Generation)}:
\begin{enumerate}
    \item L'utente invia un messaggio, insieme all'ID del documento corrente aperto
          e allo \textbf{storico completo} della conversazione.
    \item Il backend esegue una ricerca semantica e recupera i 3 documenti pi\`{u}
          rilevanti per la domanda posta.
    \item Il documento aperto viene etichettato come \texttt{[DOCUMENTO CORRENTE]},
          con i suoi metadati espliciti (firmatari, uffici, date) e fino a 3000 caratteri
          di testo estratto per contestualizzare l'LLM.
    \item Lo storico conversazione viene inviato come array di messaggi \texttt{role/content},
          permettendo al modello di rispondere in modo contestuale.
    \item L'LLM risponde in testo puro, senza Markdown (niente asterischi o elenchi puntati),
          basandosi esclusivamente sui documenti dell'archivio locale.
\end{enumerate}

La scelta del modello \texttt{GPT 5.4 Mini} e il confronto con i modelli di ragionamento
sono descritti nella sezione \textit{Alternative Valutate e Scartate} (Capitolo 2).

\subsection{Estrazione Metadati OCR tramite LLM}

Un prompt di sistema invia il testo OCR grezzo ad Azure OpenAI e ottiene
un JSON strutturato pronto per il database:

\begin{promptbox}[title={System Prompt --- Estrazione Metadati OCR da LLM}]
\textit{``Sei un assistente specializzato nell'estrazione di metadati da atti pubblici comunali
italiani. Leggi il testo OCR e restituisci ESCLUSIVAMENTE un JSON valido senza Markdown con
le chiavi: \textbf{nome} (tipologia + numero atto, es.\ Ordinanza Sindacale n.\ 12 del
05/06/2026), \textbf{descrizione} (max 2 frasi), \textbf{tipologia} (una tra le categorie
previste), \textbf{data} (YYYY-MM-DD), \textbf{data\_scadenza} (YYYY-MM-DD o null),
\textbf{uffici} (array), \textbf{firmatari} (array),
\textbf{frazioni} (array o `TUTTE').''
}
\end{promptbox}

Questo approccio converte il testo grezzo OCR in un record database completo
in un'unica chiamata API, eliminando qualsiasi inserimento manuale.

\section{Servizio Email SMTP (Gmail)}

L'invio della posta \`{e} affidato al server SMTP di Gmail (le motivazioni della scelta sono
descritte nel Capitolo 2). Il servizio gestisce:
\begin{enumerate}
    \item \textbf{Reset Password:} token monouso inviato via email, valido 15 minuti.
    \item \textbf{Notifiche di scadenza:} email HTML agli operatori a 15, 7 e 1 giorno.
\end{enumerate}

\begin{notabox}[title={Modalit\`{a} locale per le email}]
In assenza delle credenziali SMTP (Gmail), il servizio scrive il corpo dell'email
in un file \texttt{.txt} nella cartella \texttt{backend/mock\_emails/}.
\end{notabox}

\section{Riepilogo Servizi Cloud}

\begin{longtable}{>{\bfseries}p{3.5cm} p{3.5cm} p{2.5cm} p{2.8cm}}
    \toprule
    \textbf{Servizio Azure} & \textbf{Scopo} & \textbf{SDK Python} & \textbf{Fallback} \\
    \midrule
    \endhead
    Azure AI Doc Intelligence & OCR da immagini & azure-ai-formrecognizer & MockOCRService \\
    Azure Blob Storage & Archiviazione immagini & azure-storage-blob & LocalStorage \\
    Azure OpenAI Service & Chatbot RAG + Spiegazioni & openai & spaCy fallback \\
    Gmail SMTP & Invio Email & smtplib & File .txt su disco \\
    Azure Container Apps & Hosting container & Docker / az CLI & docker compose \\
    Azure Container Registry & Registro immagini & docker push & Registry locale \\
    Azure DB PostgreSQL & DB in produzione & psycopg2-binary & SQLite locale \\
    \bottomrule
    \caption{Tutti i servizi Azure utilizzati da DocuMatch}
\end{longtable}

\chapter{Pipeline di Intelligenza Artificiale}

\section{Named Entity Recognition con spaCy}

La libreria \textbf{spaCy} con il modello \texttt{it\_core\_news\_sm} identifica
automaticamente dal testo OCR:
\begin{itemize}
    \item \textbf{PER} (Person): firmatari dell'atto.
    \item \textbf{ORG} (Organization): uffici coinvolti.
    \item \textbf{LOC} (Location): frazioni o zone geografiche.
\end{itemize}

\section{Generazione Embeddings con SentenceTransformers}

Il modello \textbf{\texttt{paraphrase-multilingual-MiniLM-L12-v2}} (50+ lingue, incluso
l'italiano) converte ogni atto in un \textbf{vettore denso a 384 dimensioni}:

L'embedding viene ricalcolato: all'inserimento, durante OCR, nel bulk import,
e ad ogni modifica (PUT), mantenendo il vettore sempre sincronizzato.

\section{Cosine Similarity}

Metrica geometrica che misura l'angolo $\theta$ tra i vettori di query e documento.
Valori vicini a 1 indicano alta affinit\`{a} semantica:

\begin{equation}
    \text{sim}(A, B) = \cos\theta =
    \frac{A \cdot B}{\|A\| \cdot \|B\|} =
    \frac{\displaystyle\sum_{i=1}^{384} A_i B_i}
         {\displaystyle\sqrt{\sum_{i=1}^{384} A_i^2} \cdot
          \displaystyle\sqrt{\sum_{i=1}^{384} B_i^2}}
\end{equation}

\section{Ricerca Ibrida: Cosine + Keyword Boosting}

Per migliorare la precisione, il sistema combina la cosine similarity con un
boost morfologico sugli \textit{stem} delle parole:

\begin{itemize}
    \item \textbf{Estrazione e Stemming Custom:} La query dell'utente viene ripulita dalle \textit{stopwords} burocratiche (es. "comune", "delibera", "il", "per"). Le parole rimanenti passano attraverso una funzione di \textit{stemming} personalizzata per l'italiano (\texttt{clean\_and\_stem}) che rimuove suffissi verbali e plurali per ottenere la radice lessicale pura (es. "manutenzione" e "manutenzioni" $\rightarrow$ "manutenzi-").
    \item \textbf{Keyword Boosting (+0.05):} Se la radice estratta dalla query compare testualmente nel nome, nella descrizione o nel testo OCR del documento, il punteggio di similarit\`{a} coseno base riceve un \textit{boost} additivo (es. $+0.05$ per ogni occorrenza).
\end{itemize}

Questo approccio garantisce che i documenti in cui è presente esplicitamente la parola chiave cercata ricevano una penalizzazione minore rispetto ai documenti semanticamente affini ma che non contengono il termine esatto, unendo l'intelligenza contestuale dei transformer con la garanzia di esattezza della ricerca testuale classica.

\chapter{API REST --- Backend FastAPI}

\section{Struttura del Progetto Backend}

\section{Endpoint di Autenticazione}

\begin{longtable}{>{\ttfamily\small}p{1.5cm} >{\ttfamily\small}p{7.5cm} p{2cm} p{3cm}}
    \toprule
    \normalfont\textbf{Metodo} & \normalfont\textbf{Endpoint} &
    \normalfont\textbf{Accesso} & \normalfont\textbf{Descrizione} \\
    \midrule
    \endhead
    POST & /api/v1/auth/login & Pubblico & Login, JWT Bearer Token \\
    POST & /api/v1/auth/forgot-password & Pubblico & Richiesta reset password \\
    POST & /api/v1/auth/reset-password & Pubblico & Conferma reset con token \\
    GET  & /api/v1/auth/users & Admin Sup. & Lista operatori \\
    POST & /api/v1/auth/users & Admin Sup. & Crea nuovo operatore \\
    DELETE & /api/v1/auth/users/\{username\} & Admin Sup. & Rimuove operatore \\
    POST & /api/v1/auth/trigger-expiration-check & Auth & Trigger manuale cron scadenze \\
    GET  & /api/v1/audit/ & Admin Sup. & Registro attivit\`{a} (Audit Log) \\
    \bottomrule
    \caption{Endpoint di autenticazione e gestione utenti}
\end{longtable}

\section{Endpoint Documenti}

\begin{longtable}{>{\ttfamily\small}p{1.5cm} >{\ttfamily\small}p{7.5cm} p{2cm} p{3cm}}
    \toprule
    \normalfont\textbf{Metodo} & \normalfont\textbf{Endpoint} &
    \normalfont\textbf{Accesso} & \normalfont\textbf{Descrizione} \\
    \midrule
    \endhead
    GET  & /api/v1/documents/ & Pubblico & Lista archivio \\
    GET  & /api/v1/documents/search & Pubblico & Ricerca semantica + filtri \\
    POST & /api/v1/documents/search-by-image & Pubblico & Ricerca per foto (OCR) \\
    GET  & /api/v1/documents/\{id\} & Pubblico & Dettaglio atto \\
    GET  & /api/v1/documents/\{id\}/recommendations & Pubblico & Top 5 correlati \\
    POST & /api/v1/documents/analyze-ocr & Auth & Digitalizza copertina \\
    POST & /api/v1/documents/ & Auth & Inserimento manuale \\
    POST & /api/v1/documents/import-bulk & Auth & Bulk ZIP (JSON/CSV) \\
    POST & /api/v1/documents/export-bulk & Auth & Esporta selezione in ZIP \\
    POST & /api/v1/documents/bulk-delete & Auth & Eliminazione multipla \\
    POST & /api/v1/documents/recalculate-embeddings & Auth & Ricalcola embedding \\
    POST & /api/v1/documents/bulk-resolve-entities & Auth & Risoluzione entit\`{a} batch \\
    PUT  & /api/v1/documents/\{id\} & Auth & Modifica atto \\
    DELETE & /api/v1/documents/\{id\} & Auth & Eliminazione \\
    \bottomrule
    \caption{Endpoint di gestione documenti}
\end{longtable}

\section{Endpoint Chatbot RAG}

\begin{longtable}{>{\ttfamily\small}p{1.5cm} >{\ttfamily\small}p{7.5cm} p{2cm} p{3.2cm}}
    \toprule
    \normalfont\textbf{Metodo} & \normalfont\textbf{Endpoint} &
    \normalfont\textbf{Accesso} & \normalfont\textbf{Descrizione} \\
    \midrule
    \endhead
    POST & /api/v1/chat/ & Pubblico & Interroga l'archivio via LLM \\
    \bottomrule
    \caption{Endpoint chatbot RAG}
\end{longtable}

\section{Filtri di Ricerca e Pre-Filtering}

Il sistema implementa un approccio di \textbf{Ricerca Ibrida} che combina la precisione dei filtri deterministici (sui metadati strutturati) con la flessibilit\`{a} e la comprensione del contesto della ricerca vettoriale (sui testi destrutturati).

I filtri strutturati combinati tramite l'interfaccia utente includono:
\begin{itemize}
    \item \textbf{Ricerca Semantica}: query testuale in linguaggio naturale elaborata dal modello SentenceTransformer.
    \item \textbf{Tipologia}: selezione multipla delle categorie documentali (con logica AND/OR).
    \item \textbf{Anno}: selezione dell'anno di pubblicazione (con logica AND/OR), in sostituzione al classico e meno intuitivo range di date.
    \item \textbf{Ufficio, Firmatario, Frazione}: filtri a selezione multipla per entit\`{a} organizzative e geografiche (con logica AND/OR).
\end{itemize}

L'algoritmo opera in due step logici (approccio \textit{pre-filtering}):
\begin{enumerate}
    \item I filtri hard (esattezza) vengono applicati per primi, riducendo lo spazio di ricerca e scartando a priori i documenti che non rispettano i vincoli burocratici formali.
    \item La \textit{cosine similarity} (pertinenza) viene calcolata esclusivamente sul sottoinsieme filtrato.
\end{enumerate}
Questo meccanismo ibrido abbatte drasticamente i tempi di calcolo vettoriale e unisce l'affidabilit\`{a} di un database classico con l'intelligenza di un motore semantico.

\begin{notabox}[title={Wildcard Frazioni}]
Un documento con frazioni impostato a \texttt{"TUTTE"} viene sempre incluso
nei risultati di ricerca per zona, bypassando il filtro restrittivo.
\end{notabox}

\section{Importazione Massiva con Pandas}

L'endpoint \texttt{POST /api/v1/documents/import-bulk} accetta archivi ZIP contenenti
im\-ma\-gi\-ni e un file di metadati (\texttt{dati.json} oppure \texttt{dati.csv}).
Il parsing del CSV \`{e} delegato a \textbf{Pandas} che gestisce automaticamente separatori
ambigui, encoding diversi e valori nulli (\texttt{NaN}). Ogni documento viene indexato
(embedding calcolato) al momento dell'ingestione. Se il campo \texttt{testo\_estratto}
\`{e} vuoto o assente, il documento viene accodato per elaborazione OCR asincrona
in background, senza bloccare la risposta all'amministratore.

Le immagini estratte dallo ZIP vengono passate allo \texttt{StorageService}: se Azure
\`{e} configurato, il file viene caricato direttamente nel Blob Storage cloud; altrimenti
viene salvato localmente in modalit\`{a} Dual-Mode.

\chapter{Frontend --- Dashboard Istituzionale React}

\section{Architettura a Servizi e Viste}

Il frontend \`{e} strutturato seguendo il principio della \textbf{Separation of Concerns}
in livelli distinti:
\begin{itemize}
    \item \textbf{Componenti di presentazione} (\texttt{src/app/components/}): elementi UI
          riutilizzabili (\texttt{DocCard}, \texttt{DocFormModal}, \texttt{ChatModal}, ecc.).
    \item \textbf{Servizi} (\texttt{src/app/services/}): incapsulano le chiamate API REST
          e la logica di business, mantenendo le views snelle.
    \item \textbf{Viste} (\texttt{src/app/views/}): orchestrano la composizione e il routing,
          tenendo in stati locali i dati di pagina.
\end{itemize}

\subsection{Componenti e Viste principali}
\begin{longtable}{>{\ttfamily}p{4.5cm} p{7.5cm}}
    \toprule
    \textbf{File} & \textbf{Responsabilit\`{a}} \\
    \midrule
    \endhead
    HomeView.tsx & Home pubblica: ricerca, filtri, risultati, upload OCR \\
    AdminView.tsx & Pannello operatore: CRUD documenti, bulk import, configurazione \\
    DetailView.tsx & Dettaglio atto: testo OCR, metadati, modifica, chatbot \\
    ResetPasswordView.tsx & Pagina reset password con token da email \\
    DocCard.tsx & Scheda documento con badge affinit\`{a} AI \\
    ChatModal.tsx & Chatbot RAG con storico conversazione \\
    DocFormModal.tsx & Form inserimento / modifica con pre-compilazione OCR \\
    \bottomrule
\end{longtable}

\subsection{Servizi Frontend}

\begin{longtable}{>{\ttfamily}p{4.5cm} p{7.5cm}}
    \toprule
    \textbf{Servizio} & \textbf{Responsabilit\`{a}} \\
    \midrule
    \endhead
    authService.ts & Token JWT, login, logout, gestione utenti, trigger cron \\
    documentiService.ts & CRUD documenti, export ZIP, eliminazione bulk \\
    ricercaService.ts & Ricerca testuale (con filtri) e ricerca per immagine \\
    entitaService.ts & CRUD entit\`{a} di configurazione (uffici, tipologie, frazioni) \\
    chatService.ts & Chiamate al chatbot RAG \\
    auditService.ts & Lettura del Registro Attivit\`{a} \\
    api.ts & Interceptor HTTP: token JWT, gestione errori 401 \\
    \bottomrule
\end{longtable}

\section{Design System}

\begin{itemize}
    \item \textbf{Framework UI}: Componenti base accessibili (Radix UI) stilizzati tramite \textbf{Tailwind CSS} (stile \textit{shadcn/ui}), garantendo conformit\`{a} WCAG.
    \item \textbf{Palette e Temi}: Sistema cromatico flessibile basato su variabili CSS.
    \item \textbf{Tipografia e Iconografia}: Font di sistema sans-serif ottimizzati per la leggibilit\`{a} e set di icone coerente basato su \textbf{Lucide React}.
    \item \textbf{Responsive Design}: Layout fluido e componenti adattivi (es. \texttt{react-resizable-panels}) progettati \textit{mobile-first}.
\end{itemize}

\section{Funzionalit\`{a} UX Avanzate}

\subsection{Filtri Multi-Selezione con Logica AND/OR}

I cinque menu filtro della home page (Tipologia, Anno, Ufficio, Firmatario, Frazione)
supportano la \textbf{selezione multipla} tramite checkbox. Uno switch \textbf{TUTTI/ALMENO UNO}
per ogni dimensione configura la logica booleana. Ogni menu include una barra di ricerca
interna per filtrare le opzioni in tempo reale, utile con liste lunghe di firmatari o frazioni.
I valori vengono serializzati come array di query string e inviati all'endpoint \texttt{/search},
che li elabora sia in AND che in OR lato backend.

\subsection{Indicatore di Affinit\`{a} Semantica}

Ogni \texttt{DocCard} mostra in overlay sulla miniatura un \textbf{badge percentuale}
(cosine similarity $\times$ 100) restituito dal backend. La sezione spiegazione AI ha
altezza fissa di 140px, garantendo allineamento uniforme tra card della stessa riga.
Il prompt inviato ad Azure OpenAI richiede una sintesi istituzionale di \textbf{massimo 2-3 frasi} per le spiegazioni dei documenti correlati, al fine di mantenere le card equilibrate e leggibili senza imporre limiti di parole troppo rigidi.

\subsection{Selezione Multipla e Bulk Export}

Cliccando su uno o pi\`{u} documenti appare un menu fluttuante in basso con le azioni
\textbf{Esporta ZIP} ed \textbf{Elimina}, permettendo
operazioni in massa senza aprire ogni scheda individualmente.

\subsection{Sicurezza Token: sessionStorage vs localStorage}

Il token JWT dell'amministratore viene conservato in \textbf{sessionStorage} anziché
in \texttt{localStorage}. A differenza del \texttt{localStorage} (persistente anche
dopo la chiusura del browser), il \texttt{sessionStorage} viene azzerato alla chiusura
della tab. Questo impedisce il ricaricamento automatico del pannello admin se l'applicazione
viene riaperta, garantendo che ogni sessione parta sempre dallo stato ``non autenticato''.

\section{URL Relativi e Proxy}

Tutte le chiamate API usano URL relativi (\texttt{/api/v1/...}),
rendendo il frontend indipendente dall'indirizzo esatto del backend.
La configurazione dettagliata del proxy Nginx e il meccanismo \texttt{envsubst}
per l'iniezione dinamica di \texttt{BACKEND\_URL} sono descritti nel Capitolo 8.

\section{Autenticazione Stateless nel Frontend}

Il Bearer Token JWT viene allegato automaticamente a ogni richiesta HTTP
protetta tramite la funzione \texttt{apiFetch} nel file \texttt{api.ts}. Il gestore globale degli errori 401
provoca un logout forzato (reload della pagina), con un'eccezione esplicita
per l'endpoint \texttt{/auth/login}, che pu\`{o} legittimamente rispondere con 401 senza
provocarne il ricaricamento (la schermata di login mostra invece il messaggio di errore).

\chapter{Containerizzazione e Deploy Cloud (Azure)}

L'intera applicazione \`{e} containerizzata tramite \textbf{Docker} e deployata su
\textbf{Azure Container Apps}: ogni servizio (backend Python e frontend React) \`{e}
incapsulato in un'immagine dedicata, costruita con strategie ottimizzate per la produzione.

\section{Dockerfile Backend}

Il backend parte dall'immagine base \texttt{python:3.9-slim}, che esclude i pacchetti
non necessari riducendo la superficie d'attacco. Le fasi principali sono:

\begin{enumerate}
    \item \textbf{Dipendenze di sistema:} si installano solo \texttt{build-essential}
          e \texttt{curl} (le librerie minime per compilare alcuni pacchetti Python).
    \item \textbf{Requisiti Python:} la copia separata di \texttt{requirements.txt}
          prima del codice sorgente sfrutta la cache a layer di Docker: se le dipendenze
          non cambiano, questo step non viene rieseguito.
    \item \textbf{Pre-download modelli AI al build time:} il modello linguistico
          \texttt{it\_core\_news\_sm} (spaCy) e il modello di embedding
          \texttt{paraphrase-multilingual-MiniLM-L12-v2} (SentenceTransformer) vengono
          scaricati durante la build, non al primo avvio. Questo elimina la latenza
          iniziale in produzione (fino a 90 secondi) e rende l'immagine autonoma.
    \item \textbf{Avvio:} il server \texttt{uvicorn} viene lanciato in ascolto su
          \texttt{0.0.0.0:8000} per accettare connessioni dall'esterno del container.
\end{enumerate}

\begin{notabox}[title={Il ruolo del healthcheck Docker}]
SentenceTransformer richiede 30--90 secondi al primo avvio per caricare i pesi in
memoria. Il \texttt{start\_period: 120s} del healthcheck garantisce che il container
venga marcato come \textit{healthy} solo dopo che la pipeline NLP \`{e} completamente
operativa, evitando errori sulle prime richieste.
\end{notabox}

\section{Multi-Stage Build del Frontend}

Il Dockerfile del frontend \`{e} strutturato in \textbf{due stage} distinti per ridurre
al minimo la dimensione finale dell'immagine:

\begin{description}
    \item[\texttt{Stage 1 --- build} (Node.js 22 Alpine):] Si installa \texttt{pnpm}
          tramite \texttt{corepack}, si copiano le dipendenze e si esegue \texttt{pnpm run build}.
          Il risultato \`{e} la cartella \texttt{/dist} con l'applicazione React compilata.
          Questo stage produce un'immagine temporanea da $\sim$1.5\,GB che non
          viene mai pubblicata.
    \item[\texttt{Stage 2 --- runtime} (Nginx Alpine):] Si parte da un'immagine
          Nginx leggera ($\sim$25\,MB) e si copiano \textbf{esclusivamente} i file
          statici generati dal Stage 1 con la direttiva \texttt{COPY -{}-from=build}.
          Il risultato \`{e} un'immagine di produzione di soli $\sim$25\,MB, senza
          dipendenze di sviluppo, compilatori o codice sorgente.
\end{description}

\section{Nginx con envsubst}

Il container frontend usa \textbf{Nginx Alpine} come web server. Il problema classico
del frontend containerizzato \`{e} che l'URL del backend non \`{e} noto al momento della
\textit{build}: cambia tra sviluppo locale, staging e produzione Azure.

La soluzione adottata \`{e} il meccanismo \textbf{envsubst} (dalla libreria \texttt{gettext}):
\begin{enumerate}
    \item Il file \texttt{nginx.conf} usa il \textbf{segnaposto} \texttt{\$\{BACKEND\_URL\}}
          invece di un indirizzo hardcoded.
    \item Al momento dell'avvio del container, il \texttt{CMD} del Dockerfile
          esegue \texttt{envsubst} che sostituisce \texttt{\$\{BACKEND\_URL\}} con il valore
          reale della variabile d'ambiente, generando il file di configurazione definitivo.
    \item Nginx viene poi avviato con il file gi\`{a} compilato.
\end{enumerate}

Il template \texttt{nginx.conf} definisce tre \texttt{location} block distinti:
\begin{description}
    \item[\texttt{/api/}:] tutte le richieste API vengono inoltrate al backend FastAPI
          tramite \texttt{proxy\_pass} con timeout esteso a 120s
          (per le operazioni OCR che possono durare fino a 30--90 secondi).
    \item[\texttt{/static/}:] i file statici (immagini in modalit\`{a} locale)
          vengono proxati al backend; in produzione Azure vengono serviti direttamente
          tramite URL del Blob Storage.
    \item[\texttt{/}:] tutti gli altri percorsi servono \texttt{index.html} con
          \texttt{try\_files \$uri \$uri/ /index.html}, abilitando il routing lato client
          di React Router senza 404 su navigazione diretta.
\end{description}

\begin{notabox}[title={Valore di default per sviluppo locale}]
Il \texttt{CMD} del Dockerfile usa la sintassi \texttt{\$\{BACKEND\_URL:-http://backend:8000\}}:
se la variabile non \`{e} impostata, il fallback \`{e} \texttt{http://backend:8000},
che corrisponde al nome del servizio backend in Docker Compose.
In produzione Azure, \texttt{BACKEND\_URL} viene impostato dall'esterno nel pannello
\textit{Environment Variables} del Container App.
\end{notabox}

\section{Deploy su Azure Container Apps}

Il deploy avviene in quattro passi principali:

\begin{enumerate}
    \item \textbf{Provisioning delle risorse Azure:} creazione dell'Azure Container Registry
          (\texttt{acrdocumatch}), del database PostgreSQL Flexible Server con il DB
          \texttt{documatch\_db}, dell'Azure Blob Storage per le immagini, del servizio
          Azure AI Document Intelligence (OCR) e dell'Azure OpenAI Service con il
          deployment del modello \texttt{GPT 5.4 Mini}.
    \item \textbf{Build e Push:} le immagini Docker vengono costruite localmente
          (o tramite pipeline CI/CD) e caricate nell'Azure Container Registry
          con \texttt{az acr build} o \texttt{docker push}.
    \item \textbf{Creazione dei Container Apps:} il backend viene deployato su porta
          8000 con \textit{Ingress External} e le variabili d'ambiente:
          \texttt{DATABASE\_URL}, \texttt{SECRET\_KEY},
          \texttt{AZURE\_DOCUMENT\_INTELLIGENCE\_KEY/ENDPOINT},
          \texttt{AZURE\_STORAGE\_CONNECTION\_STRING},
          \texttt{AZURE\_OPENAI\_KEY/ENDPOINT}. \\
          Il frontend viene deployato su porta 80 con \texttt{BACKEND\_URL} impostato
          all'URL del Container App del backend.
    \item \textbf{Configurazione CORS:} la variabile \texttt{CORS\_ALLOWED\_ORIGINS}
          viene impostata all'URL esatto del frontend. Con \texttt{ENVIRONMENT=production}
          e CORS wildcard (\texttt{*}), il backend si rifiuta di avviarsi per prevenire
          accessi non autorizzati da origini arbitrarie.
\end{enumerate}

\chapter{Sicurezza e Autenticazione}

\section{Token JWT (HS256 / HMAC-SHA256) Stateless}

Autenticazione stateless senza sessioni server-side. Viene usata la libreria
\textbf{PyJWT} con algoritmo \textbf{HS256} per generare token JWT standard
(Header.Payload.Signature) firmati con la \texttt{SECRET\_KEY}:

\section{RBAC: API Pubbliche e Protette in Scrittura}
Il framework implementa un sistema di controllo degli accessi (Role-Based Access Control) granulare a livello di endpoint:
\begin{itemize}
    \item \textbf{Lettura pubblica:} Trattandosi di un simulatore di albo pretorio, le operazioni di lettura (GET) e di ricerca semantica (POST per image-search) non richiedono alcun token JWT, rendendo i documenti finali accessibili a tutti i cittadini senza autenticazione.
    \item \textbf{Scrittura protetta:} Qualsiasi endpoint che implichi una modifica dei dati o dei modelli (creazione, aggiornamento, eliminazione documenti, generazione OCR e ricalcolo embeddings) \`e rigidamente protetto lato server dalla dipendenza \texttt{Depends(get\_current\_admin)}. In assenza di un header \texttt{Authorization: Bearer <token>} valido, FastAPI respinge intrinsecamente la chiamata con errore HTTP 401 Unauthorized, rendendo impossibile l'alterazione del database eludendo la UI del frontend.
\end{itemize}

\section{Password Hashing con Bcrypt}

Bcrypt implementa \textit{key stretching}: ogni hash richiede un tempo di calcolo
parametrizzabile, rendendo il brute-force da GPU impraticabile.

\section{Reset Password con Token Monouso}

\begin{enumerate}
    \item L'operatore richiede il reset con il proprio username.
    \item Il backend genera un token casuale (128 bit, hex) con scadenza 15 minuti.
    \item Il token viene salvato nel DB. Un link viene inviato via email.
    \item L'operatore clicca il link; il backend verifica il token e aggiorna la password
          (bcrypt). Il token viene invalidato impostandolo a \texttt{null}.
\end{enumerate}

\section{Protezione Anti-ZipBomb}
L'endpoint di caricamento massivo di documenti tramite archivi ZIP (\texttt{/api/v1/documents/import-bulk}) \`e vulnerabile ad attacchi di tipo \textit{Zip Bomb} o \textit{Decompression Bomb}, dove un piccolo archivio (es. 42 KB) si decomprime in Petabyte di dati riempiendo la memoria del server. 
Per mitigare questa vulnerabilit\`a, il codice ispeziona i metadati dell'archivio (tramite \texttt{z.infolist()}) prima di effettuare la lettura. Vengono imposti i seguenti limiti rigidi:
\begin{itemize}
    \item \texttt{MAX\_ZIP\_FILES}: massimo 1000 file nell'archivio.
    \item \texttt{MAX\_UNCOMPRESSED\_SIZE}: dimensione massima espansa bloccata a 100 MB.
\end{itemize}
Se questi limiti vengono superati, l'estrazione \`e annullata e la richiesta respinta con errore 400 (Bad Request).

\section{Sistema di Accountability (Audit Log)}
Un sistema documentale richiede la tracciabilit\`a di tutte le azioni sensibili per la sicurezza aziendale.
\`E stato implementato un \texttt{AuditLog} su database relazionale:
\begin{itemize}
    \item Il modello traccia \texttt{timestamp}, \texttt{username}, \texttt{action} (es. CREATE\_DOC, UPDATE\_DOC, MASS\_IMPORT) e dettagli esplicativi.
    \item Ogni endpoint REST che esegue modifiche salva a posteriori il log dell'operazione.
    \item Nell'interfaccia React \`e stata aggiunta un'area amministrativa ``Registro Attivit\`a'' visibile solo agli amministratori autorizzati (Supreme Admin) per un'analisi real-time delle scritture.
\end{itemize}

\chapter{Sistema di Notifiche Schedulate}

\section{APScheduler in Background}

Lo scheduler viene avviato nel lifecycle \texttt{lifespan} di FastAPI tramite
\textbf{APScheduler} con un cron job configurato per le ore \textbf{08:00 ogni giorno}.
La logica di notifica implementa un sistema di \textbf{escalation graduale}:
\begin{itemize}
    \item \textbf{15 giorni prima:} prima notifica email a tutti gli operatori registrati.
    \item \textbf{7 giorni prima:} secondo promemoria.
    \item \textbf{1 giorno prima:} notifica urgente di imminente scadenza.
\end{itemize}
Per evitare spam, il sistema controlla che la notifica per un determinato documento
non sia gi\`{a} stata inviata nella stessa finestra temporale.

L'endpoint \texttt{POST /api/v1/documents/trigger-expiry-check} consente
all'amministratore di forzare il controllo manualmente dalla UI (pulsante ``Trigger Email Scadenze''
nella scheda Impostazioni del pannello admin), utile per test e demo.
In modalit\`{a} demo, il trigger considera qualsiasi documento in scadenza
nei prossimi 15 giorni (non solo quelli esattamente ai tre giorni soglia),
inviando immediatamente le notifiche per tutti i documenti in procinto di scadere.

\chapter{Integrazione dell'IA Generativa --- Dichiarazione di Trasparenza}

\section{Strumento Utilizzato}

In conformit\`{a} con le indicazioni del docente, si dichiara che per lo sviluppo del
progetto \`{e} stato utilizzato lo strumento di AI Generativa \textbf{Antigravity IDE}
(basato su Google Gemini e Anthropic Claude Sonnet), integrato in Visual Studio Code.

L'IA \`{e} stata usata come \textbf{strumento di accelerazione}, non come sostituto
della progettazione. Ogni componente \`{e} stato compreso, rivisto e adattato.

\section{Distribuzione del Lavoro}

\begin{longtable}{>{\bfseries}p{3.2cm} p{4.5cm} p{4.5cm}}
    \toprule
    \textbf{Area} & \textbf{Contributo AI} & \textbf{Contributo Studente} \\
    \midrule
    \endhead
    Architettura & Pattern Dual-Mode, SOLID/DIP & Scelte Azure, adattamento dominio \\
    Pipeline NLP & Bozza spaCy + SentenceTransformers & Tuning, integrazione, hybrid boost \\
    Frontend React & Bozza componenti e hook TypeScript & UX istituzionale, CSS custom \\
    Backend FastAPI & Struttura router, DI & Logica business, filtri, migrazioni \\
    Scheduler & Configurazione APScheduler & Dual-mode demo/cron, lifespan \\
    Deploy Azure & Pattern envsubst Nginx & Test, verifica, adattamento \\
    Sicurezza & HMAC token, bcrypt & Reset flow, CORS produzione \\
    Relazione & Struttura, formattazione & Contenuto tecnico, revisione \\
    \bottomrule
    \caption{Distribuzione lavoro studente / AI}
\end{longtable}

\section{Prompt Principali Utilizzati}

Di seguito vengono riportati i prompt reali estratti dalle conversazioni esportate con l'Intelligenza Artificiale, a dimostrazione dell'esteso e complesso flusso di lavoro svolto per la realizzazione di DocuMatch.

\subsection{Fase 1: UI/UX e Frontend React}

\begin{promptbox}[title={Prompt 1 --- Progettazione Interfaccia Base}]
\textbf{Utente:} \textit{`Crea un sito web per la gestione delle schede documentali del comune di Montalto Uffugo. La piattaforma si chiamerà DocuMatch e deve avere una home page con tutte le schede, una barra di ricerca con filtri combinabili e una modalità di ricerca OCR da immagine tramite uno switch dedicato. L'area admin deve permettere la modifica delle schede.''}

\textbf{Risposta e Azioni:} L'AI ha generato l'architettura React in \texttt{Documatch/src/app}, implementando il layout grid a 3 colonne, la barra di ricerca full-text e i filtri combinabili.
\end{promptbox}

\begin{promptbox}[title={Prompt 2 --- Ricerca OCR per Immagine}]
\textbf{Utente:} \textit{`Aggiungi una sezione di ricerca tramite OCR con un'area di caricamento drag-and-drop ben visibile. Sostituisci lo switch poco evidente con due tab distinte in cima alla barra di ricerca: una per la ricerca testuale e una per la ricerca per immagine.''}

\textbf{Risposta e Azioni:} L'AI ha riprogettato la barra di ricerca aggiungendo due tab visibili in cima (Ricerca testuale / Ricerca per immagine). È stata implementata una dropzone drag-and-drop per i formati JPG/PNG/PDF.
\end{promptbox}

\begin{promptbox}[title={Prompt 3 --- Integrazione OCR nel Form Admin}]
\textbf{Utente:} \textit{`Aggiungi l'upload OCR simulato direttamente nel form di inserimento.''}

\textbf{Risposta e Azioni:} Modificato il \texttt{DocFormModal.tsx}. L'AI ha aggiunto una sezione di caricamento file in cima al form che, al drop dell'immagine, avvia un'analisi OCR simulata e pre-compila automaticamente i campi Numero, Nome, Descrizione e Tipologia.
\end{promptbox}

\begin{promptbox}[title={Prompt 4 --- Configurazione Dinamica Metadati}]
\textbf{Utente:} \textit{`Trasforma i campi a testo libero (Uffici, Firmatari) in valori preimpostati e modificabili dall'amministratore. Aggiungi una sezione ``Configurazione'' nell'area admin che permetta di gestire queste entità tramite interfaccia grafica anziché inserimento manuale.''}

\textbf{Risposta e Azioni:} Aggiunta la sezione ``Configurazione'' nell'area admin. L'AI ha trasformato i campi di testo libero (Uffici, Firmatari) in griglie di pill/checkbox multi-select estrapolate dal database, migliorando la consistenza dei dati.
\end{promptbox}

\subsection{Fase 2: Backend e Integrazione API}

\begin{promptbox}[title={Prompt 5 --- Allineamento Tipi Backend/Frontend}]
\textbf{Utente:} \textit{`Verifica che il backend sia correttamente allineato con il frontend di Documatch. Controlla che gli schemi Pydantic e le interfacce TypeScript siano consistenti per tutti i campi, le date e le entità collegate, e correggi ogni discrepanza trovata.''}

\textbf{Risposta e Azioni:} L'AI ha scansionato \texttt{Documatch/src/app/types} e \texttt{backend/app/schemas/}. Notando discrepanze sui formati di date e sulle entità collegate, ha unificato le dichiarazioni TypeScript e Pydantic.
\end{promptbox}

\begin{promptbox}[title={Prompt 6 --- Semplificazione Schemi e Documentazione Codice}]
\textbf{Utente:} \textit{`Semplifica gli schemi Pydantic dove possibile ed aggiungi commenti esplicativi diffusi in tutto il codice. Il sistema deve essere comprensibile.''}

\textbf{Risposta e Azioni:} L'AI ha analizzato l'albero dei modelli. Anzich\`{e} rimuovere funzionalit\`{a}, ha unificato e ripulito ulteriormente gli schemi Pydantic e le dichiarazioni, commentando ampiamente il codice.
\end{promptbox}

\subsection{Fase 3: Intelligenza Artificiale e Cloud}

\begin{promptbox}[title={Prompt 7 --- Porting su Cloud Azure}]
\textbf{Utente:} \textit{`Implementa il porting completo del progetto su Microsoft Azure. Adatta il backend per usare PostgreSQL tramite variabile d'ambiente, sostituisci i mock locali con gli SDK Azure reali per Blob Storage e Document Intelligence, e configura i container per il deploy su Azure Container Apps.''}

\textbf{Risposta e Azioni:} L'AI ha tracciato la roadmap infrastrutturale: adattamento di \texttt{database.py} per PostgreSQL via URL d'ambiente, sostituzione dei Mock locali con SDK Azure reali (Blob Storage) e configurazione dei container. Le modifiche sono state implementate in \texttt{backend/app/core/}.
\end{promptbox}

\begin{promptbox}[title={Prompt 8 --- Chatbot RAG per interrogazione Atti}]
\textbf{Utente:} \textit{`Implementa un chatbot interattivo stile AI generativa che permetta all'utente di interrogare direttamente i documenti dell'archivio. Usa Azure OpenAI con pattern RAG (Retrieval-Augmented Generation): cerca i documenti più rilevanti con la ricerca vettoriale e iniettali nel contesto del prompt per ottenere risposte basate esclusivamente sull'archivio locale.''}

\textbf{Risposta e Azioni:} L'AI ha esteso \texttt{generative\_ai.py} e integrato l'API di Azure OpenAI per implementare un Chatbot Documentale RAG che risponde in modo semantico basandosi esclusivamente sul contenuto del documento corrente o sui metadati estratti.
\end{promptbox}

\begin{promptbox}[title={Prompt 9 --- Gestione Campi Vuoti nell'Estrazione OCR}]
\textbf{Utente:} \textit{`Gestisci correttamente i campi non trovati dall'OCR: i campi stringa non rilevati devono essere impostati a \texttt{null}, le liste a array vuoto. Non inserire mai etichette di default come ``Altro'' o ``Sconosciuto'' quando l'informazione non è presente nel testo.''}

\textbf{Risposta e Azioni:} L'AI ha adeguato i servizi di validazione in \texttt{nlp.py} e nel prompt di Azure OpenAI, assicurandosi che i campi non estratti venissero impostati correttamente a \texttt{null} o lista vuota, senza inserire etichette errate.
\end{promptbox}

\subsection{Fase 4: DevOps, Debugging e Consegna}

\begin{promptbox}[title={Prompt 10 --- Correzione Healthcheck e Avvio Docker}]
\textbf{Utente:} \textit{`Risolvi l'errore di avvio del container backend (\texttt{dependency failed to start}). Il backend impiega troppo ad avviarsi a causa del caricamento dei modelli AI: aumenta il \texttt{start\_period} dell'healthcheck in \texttt{docker-compose.yml} e correggi le policy \texttt{depends\_on} affinché il frontend attenda la piena operatività del backend prima di ricevere traffico.''}

\textbf{Risposta e Azioni:} Verificando i log, l'AI ha individuato che il caricamento dei modelli AI locali causava il timeout dell'healthcheck di Docker. Ha corretto \texttt{docker-compose.yml} aumentando gli \texttt{start\_period} e sistemato le policy di \texttt{depend\_on} tra frontend e backend.
\end{promptbox}

\vspace{1em}

\begin{promptbox}[title={Prompt 11 --- Code Audit e Pulizia della Codebase}]
\textbf{Utente:} \textit{`Fai un controllo sul progetto. Controlla bene se c'è del codice morto e rimuovilo, controlla la pulizia generale.''}

\textbf{Risposta e Azioni:} L'AI ha eseguito uno scan completo. Sono stati rimossi script di test obsoleti, ripuliti gli import dismessi in \texttt{main.py} e \texttt{App.tsx}, e rimossi mockup temporanei superati dall'effettiva implementazione API, ottenendo una repository pulita per la consegna.
\end{promptbox}

\subsection{Fase 5: UX Avanzata e Correzioni Frontend}

\begin{promptbox}[title={Prompt 12 --- Correzione Credenziali Errate Ricaricano la Pagina}]
\textbf{Utente:} \textit{`Ogni volta che sbaglio le credenziali la pagina si aggiorna e mi riporta alla home.''}

\textbf{Risposta e Azioni:} L'AI ha individuato il bug in \texttt{api.ts}: il gestore globale degli errori 401 ricaricava la pagina anche durante il login, prima che il messaggio di errore potesse essere mostrato. È stata aggiunta un'eccezione esplicita per l'endpoint \texttt{/auth/login}, consentendo al componente di catturare l'errore e mostrare la scritta rossa senza redirect.
\end{promptbox}

\begin{promptbox}[title={Prompt 13 --- Rimozione Completa del Campo Numero Protocollo}]
\textbf{Utente:} \textit{`Rimuovi completamente il campo numero dall'intero sistema.''}

\textbf{Risposta e Azioni:} L'AI ha rimosso la colonna \texttt{numero} dal modello SQLAlchemy, dagli schemi Pydantic, dal form React (\texttt{DocFormModal.tsx}), dalle \texttt{DocCard}, dalla tabella admin e dai dati mock dell'OCR, mantenendo la coerenza full-stack senza lasciare dati orfani nel database.
\end{promptbox}

\begin{promptbox}[title={Prompt 14 --- Stabilizzazione Visiva della Barra di Scorrimento}]
\textbf{Utente:} \textit{`Si ha un effetto visivo fastidioso quando appare o scompare la barra laterale di scorrimento.''}

\textbf{Risposta e Azioni:} Quando il numero di documenti filtrati non generava overflow, il layout si ``scattava'' lateralmente. L'AI ha aggiunto la regola CSS \texttt{overflow-y: scroll} sul tag \texttt{html} per riservare permanentemente lo spazio della scrollbar, eliminando ogni spostamento visivo al cambio dei risultati.
\end{promptbox}

\begin{promptbox}[title={Prompt 15 --- Riorganizzazione Area Admin con DocCard Integrate}]
\textbf{Utente:} \textit{`Nell'area admin mostra solo la possibilità di aggiunta. La modifica e l'eliminazione devono potersi fare direttamente dalla scheda, non dalla lista.''}

\textbf{Risposta e Azioni:} L'AI ha rimosso la tabella-lista da \texttt{AdminView.tsx}, trasformandola in un pannello compatto con pulsanti (Aggiungi / Importa Massivo / Configurazione). I bottoni Modifica ed Elimina sono stati inseriti direttamente in \texttt{DocCard.tsx} e \texttt{DetailView.tsx}, visibili solo quando l'utente \`{e} autenticato come amministratore.
\end{promptbox}

\begin{promptbox}[title={Prompt 16 --- Barra di Ricerca nel Form e Selezione Multipla Frazioni}]
\textbf{Utente:} \textit{`Nel form di inserimento aggiungi la ricerca per filtrare le opzioni, e la possibilit\`{a} di scegliere pi\`{u} frazioni.''}

\textbf{Risposta e Azioni:} L'AI ha esteso il modello \texttt{Document} trasformando il campo \texttt{frazione} (stringa singola) in un array JSON serializzato, creato una migrazione additive per SQLite e aggiornato \texttt{DocFormModal.tsx} con un campo di ricerca che filtra in tempo reale le pill per Tipologia, Uffici, Firmatari e Frazioni.
\end{promptbox}

\begin{promptbox}[title={Prompt 17 --- Logica Wildcard ``Tutte le Frazioni''}]
\textbf{Utente:} \textit{`Implementa la logica per cui se viene selezionata una frazione devono essere considerati anche i documenti validi per tutte le frazioni.'}

\textbf{Risposta e Azioni:} L'AI ha aggiornato l'endpoint \texttt{/search} in \texttt{documents.py}: quando viene applicato il filtro su una frazione specifica, vengono ora inclusi automaticamente anche i documenti che hanno \texttt{frazioni} impostato a \texttt{"TUTTE"}, trattandolo come wildcard istituzionale a copertura dell'intero territorio comunale.
\end{promptbox}

\subsection{Fase 6: Deploy Cloud e Ottimizzazioni Azure}

\begin{promptbox}[title={Prompt 18 --- Storage Blob per Importazione Massiva}]
\textbf{Utente:} \textit{`Nell'importazione massiva le immagini non vengono salvate su Azure Blob ma solo in locale.''}

\textbf{Risposta e Azioni:} L'AI ha riscritto il gestore ZIP in \texttt{documents.py}: ogni immagine estratta dall'archivio viene ora passata al \texttt{StorageService}, che la carica su Azure Blob Storage (o la salva localmente in Dual-Mode). \`{E} stata aggiornata anche l'esportazione ZIP per recuperare le immagini correttamente sia dal blob cloud che dal filesystem locale.
\end{promptbox}

\begin{promptbox}[title={Prompt 19 --- Chatbot con Memoria, Contesto e Senza Markdown}]
\textbf{Utente:} \textit{`Il chatbot risponde male, inserisce asterischi, non ricorda la conversazione e confonde il documento corrente con gli altri. Specifica nei prompt di non utilizzare markdown e il documento corrente deve essere sempre ben distinto dagli altri.''}

\textbf{Risposta e Azioni:} L'AI ha rielaborato il prompt di sistema in \texttt{chat.py} con direttive rigide (niente Markdown, nessun asterisco) e ha distinto il documento corrente dagli altri usando l'etichetta \texttt{[DOCUMENTO CORRENTE ATTUALMENTE APERTO]}. Il frontend \`{e} stato corretto per inviare lo storico completo della conversazione ad ogni chiamata, abilitando la memoria della chat.
\end{promptbox}

\begin{promptbox}[title={Prompt 20 --- Correzione Estrazione Titolo OCR tramite LLM}]
\textbf{Utente:} \textit{`Il titolo estratto tramite OCR \`{e} sbagliato: scrive la descrizione invece del numero e tipo dell'atto.''}

\textbf{Risposta e Azioni:} Il cambio di modello OpenAI aveva alterato l'interpretazione del campo \texttt{nome}. L'AI ha reso il prompt pi\`{u} prescrittivo: \textit{``Includi SEMPRE il tipo e il numero dell'atto se presenti (es.\ `Ordinanza Sindacale n.~12'). Evita descrizioni generiche.''}, separando nettamente il titolo identificativo dalla descrizione testuale.
\end{promptbox}

\begin{promptbox}[title={Prompt 21 --- Risoluzione Bug Email Notifiche di Scadenza}]
\textbf{Utente:} \textit{`Il sistema di email per le notifiche di scadenza non funziona.''}

\textbf{Risposta e Azioni:} Il pulsante ``Trigger Email Scadenze'' usava una funzione \texttt{apiFetch} non importata nel componente React. L'AI ha spostato la chiamata nel corretto \texttt{authService.ts} e ha migliorato la logica backend: il cron job invia avvisi a 15, 7 e 1 giorno dalla scadenza (con gradazione di urgenza), mentre la modalit\`{a} demo include tutti gli atti in scadenza entro 15 giorni.
\end{promptbox}

\begin{promptbox}[title={Prompt 22 --- Visibilit\`{a} Password e Ancoraggio Spiegazione AI}]
\textbf{Utente:} \textit{`Aggiungi un'icona occhio in tutti i campi password, e ancora la descrizione di similarit\`{a} in basso nelle card.''}

\textbf{Risposta e Azioni:} L'AI ha aggiunto il toggle di visibilit\`{a} password (icona \texttt{Eye}/\texttt{EyeOff} da \texttt{lucide-react}) nei form di login e reset. In \texttt{DocCard.tsx} ha usato il layout Flexbox con \texttt{h-full} e \texttt{flex-1} per spingere il pannello della spiegazione AI sempre al fondo della card, garantendo allineamento uniforme in griglia multi-colonna.
\end{promptbox}

\begin{promptbox}[title={Prompt 23 --- Allineamento Tipologie Esistenti nel Database}]
\textbf{Utente:} \textit{`Allinea nel database le tipologie dei documenti gi\`{a} esistenti per renderle consistenti con quelle configurate nel sistema.''}

\textbf{Risposta e Azioni:} L'AI ha eseguito uno script di migrazione dati che ha aggiornato 15 record con le vecchie diciture abbreviate (``Determina'', ``Ordinanza'', ``Atto'') convertendole nelle categorie estese ufficiali (``Determinazione Dirigenziale'', ``Ordinanza Sindacale'', ``Altro''), garantendo la consistenza con i filtri dell'interfaccia.
\end{promptbox}

\begin{promptbox}[title={Prompt 24 --- Protezione Anti-Zip Bomb per Importazione Massiva}]
\textbf{Utente:} \textit{`Aggiungi una protezione contro possibili Zip Bomb nell'endpoint di importazione massiva.''}

\textbf{Risposta e Azioni:} L'AI ha aggiunto in \texttt{documents.py} un controllo preventivo sui metadati dell'archivio ZIP prima dell'estrazione: vengono imposti il limite di 1000 file e 100 MB di dimensione decompressa massima. Se superati, l'operazione viene annullata con errore HTTP 400, prevenendo attacchi di tipo Decompression Bomb.
\end{promptbox}

\chapter{Conclusioni e Sviluppi Futuri}

\section{Risultati Raggiunti}

Il progetto \textbf{DocuMatch} ha centrato pienamente tutti gli obiettivi accademici
e tecnologici della traccia di Sistemi Distribuiti e Cloud Computing.

L'adozione di \textbf{Azure Container Apps} al posto delle classiche VM IaaS ha dimostrato
concretamente i vantaggi del paradigma PaaS: auto-scaling, networking e TLS automatici,
zero gestione del sistema operativo.

La \textbf{pipeline AI} (spaCy NER, SentenceTransformers, Azure OpenAI,
Azure AI Document Intelligence) ha trasformato un semplice archivio digitale in uno
strumento documentale intelligente, con una ricerca semantica che supera ampiamente
i sistemi keyword-based tradizionali.

Il pattern \textbf{Dual-Mode} ha consentito uno sviluppo efficiente senza sprecare
crediti Azure, garantendo al contempo la piena funzionalit\`{a} in produzione cloud.

\section{Competenze Acquisite}

\begin{itemize}
    \item Progettazione di architetture cloud-native (IaaS, PaaS, SaaS).
    \item Containerizzazione Docker e ottimizzazione con Multi-Stage Build.
    \item Utilizzo degli SDK Azure (Blob Storage, Document Intelligence, OpenAI).
    \item Implementazione di algoritmi di similarit\`{a} semantica su embedding vettoriali.
    \item Sviluppo di API REST ASGI con FastAPI e SQLAlchemy ORM.
    \item Utilizzo documentato e critico dell'IA Generativa nello sviluppo software.
\end{itemize}

\section{Sviluppi Futuri}

\begin{itemize}
    \item \textbf{Supporto PDF Nativo:} integrare \texttt{PyPDF2} o \texttt{pdfminer} per PDF multipagina.
    \item \textbf{Ricerca con pgvector:} sostituire il calcolo in-memory con l'estensione
          PostgreSQL \texttt{pgvector} per ricerca ANN a livello database.
    \item \textbf{Autenticazione SPID / CIE:} identit\`{a} digitale federata ministeriale.
    \item \textbf{Fine-Tuning dei Modelli:} corpus di atti burocratici italiani per
          massimizzare la precisione nel gergo tecnico-legale.
    \item \textbf{Indicizzazione Incrementale:} aggiornare solo gli embedding modificati.
\end{itemize}

\end{document}
